\chapter{Periodic phases and the long-period transition limit}\label{sec:examples-phase-periodic}
The Hamiltonian in Lemma \ref{lem:phase-hamiltonian} describes many solutions outside the finite-energy class. We develop the infinite-energy examples suggested in \cite{ChAC}, Section 2.2 and Exercise 2.2, and unpack the periodic orbit appearing in Lemma \ref{lem:phase-stable-ode}. Throughout this section $\varepsilon=1$ unless explicitly rescaled.

\begin{proposition}[Every bounded entire profile stays between the wells]\label{prop:phase-bounded-ode-range}
If $f\in C^2(\mathbb R)$ is bounded and solves $f''=f^3-f$, then $|f|\leq1$. If it takes the value $+1$ or $-1$, it is that constant. For its Hamiltonian $C=(f')^2/2-W(f)$ one has $-1/4\leq C\leq0$.
\end{proposition}
\begin{proof}
Suppose $f(t_0)=b>1$. Reverse time if necessary so that $f'(t_0)\geq0$. As long as $f\geq b$, its second derivative is at least $b^3-b>0$. Integration shows that it remains above $b$ with increasing derivative for every later time, contradicting boundedness. Apply the same argument to $-f$ to exclude values below $-1$. At a well, a nonzero derivative would cross outside the allowed range; with zero derivative, ODE uniqueness gives the constant solution.

Since $W(f)\leq1/4$, the Hamiltonian is at least $-1/4$. If it were positive, the identity $(f')^2=2(W(f)+C)\geq2C$ would make $f'$ nonvanishing with constant sign and magnitude bounded below on the line. That again contradicts boundedness. This proves both endpoints of the permitted range without assuming that a global maximum is attained.
\end{proof}

\begin{theorem}[Classification of bounded entire profiles]\label{thm:phase-bounded-ode-classification}
Every bounded entire solution is a constant in $\{-1,0,1\}$, a translated and possibly reversed heteroclinic, or a periodic profile with amplitude $a\in(0,1)$ oscillating between $-a$ and $a$. For each amplitude there is, up to translation and reversal, exactly one such periodic solution. It satisfies
\[
 (f')^2=\frac12(a^2-f^2)(2-a^2-f^2),\qquad C=-W(a).
\]
The periodic profiles and the zero constant have infinite total energy on the line. This extends the finite-energy classification of Theorem \ref{thm:phase-ode-classification}, rather than identifying these distinct classes, and supplies the bounded infinite-energy examples asked for in \cite{ChAC}, Exercise 2.2.
\end{theorem}
\begin{proof}
Proposition \ref{prop:phase-bounded-ode-range} treats the wells. At $C=0$, separation on $(-1,1)$ gives precisely the heteroclinics already classified. At $C=-1/4$, the identity forces $f'=0$ and $W(f)=1/4$, so $f=0$. For $-1/4<C<0$, there is a unique $a\in(0,1)$ with $W(a)=-C$, since $W$ is strictly decreasing on $(0,1)$. Factoring $2(W(f)-W(a))$ gives the displayed expression. On the permitted range $(-1,1)$ this is nonnegative exactly when $|f|\leq a$.

The separated time integral from $-a$ to $a$ is finite: both endpoint zeros are simple. Starting at $(f,f')=(0,\sqrt{2(W(0)-W(a))})$, the solution reaches $a$, turns because $f''=a^3-a<0$, reaches $-a$ and turns again. Reflection at the extrema and ODE uniqueness produce a smooth periodic solution. The conserved orbit contains every possible nonconstant datum at this energy, proving uniqueness up to its phase and reversal. Its energy over a period is positive, so summing disjoint periods gives infinite total energy. The constant zero has the positive constant energy density $1/4$.
\end{proof}

\begin{proposition}[The fundamental period is strictly increasing]\label{prop:phase-period-function}
For amplitude $a\in(0,1)$ the least period is
\[
 T(a)=4\sqrt2\int_0^{\pi/2}
       \frac{d\theta}{\sqrt{\,2-a^2(1+\sin^2\theta)\,}}.
\]
It is smooth and strictly increasing, with
$\lim_{a\downarrow0}T(a)=2\pi$ and $\lim_{a\uparrow1}T(a)=+\infty$.
For $a\geq1/2$ it satisfies
$T(a)\leq C(1+|\log(1-a^2)|)$.
\end{proposition}
\begin{proof}
One quarter-period is the time from $0$ to $a$. In its separated integral substitute $f=a\sin\theta$; the factor $a\cos\theta$ cancels and yields the formula. This orbit is strictly monotone between successive extrema, so the resulting full return is the least period. On any compact amplitude subinterval the denominator stays positive. Differentiating under the integral gives
\[
 T'(a)=4\sqrt2\,a\int_0^{\pi/2}
 \frac{1+\sin^2\theta}{[2-a^2(1+\sin^2\theta)]^{3/2}}\,d\theta>0.
\]
Dominated convergence at zero gives $2\pi$. At one, the integrand increases to $1/\cos\theta$, whose integral diverges, proving the other limit by monotone convergence.

Write the squared denominator as $2(1-a^2)+a^2\cos^2\theta$. For $a\geq1/2$ and $r=\pi/2-\theta$, use $\cos\theta=\sin r\geq2r/\pi$. The integral is bounded by a constant times $\int_0^{\pi/2}((1-a^2)+r^2)^{-1/2}dr$. Its elementary antiderivative gives the logarithmic upper bound.
\end{proof}

\begin{example}[Small oscillations and a single long-period layer]\label{ex:phase-periodic-limits}
Normalize the increasing zero of the periodic profile by $f_a(0)=0$ and $f_a'(0)>0$. Then
\[
 a^{-1}f_a(t)\longrightarrow\sin t\quad(a\downarrow0),\qquad
 f_a(t)\longrightarrow q(t)\quad(a\uparrow1)
\]
in $C^\infty$ on every compact interval. Each long-period solution contains both an increasing and a decreasing transition; the second moves outside every fixed interval under this normalization.
\end{example}
\begin{proof}
The initial derivative is $f_a'(0)=a\sqrt{1-a^2/2}$. Thus $v_a=f_a/a$ solves
$v_a''=a^2v_a^3-v_a$ with $v_a(0)=0$ and $v_a'(0)=\sqrt{1-a^2/2}$. Since $|v_a|\leq1$, the integral form of this ODE and Gronwall's inequality give convergence with first derivatives to the solution of $v''=-v$, $v(0)=0$, $v'(0)=1$. The equation and successive differentiations give every higher derivative. This limit is $\sin t$.

At the other endpoint the initial data converge to $(0,1/\sqrt2)$ for the unchanged equation $f''=f^3-f$. The same continuous-dependence argument, using $|f_a|\leq1$, yields the heteroclinic $q$ on compact intervals. Its decreasing zero occurs at $T(a)/2$, which diverges by Proposition \ref{prop:phase-period-function}. This explains why local convergence to one layer is compatible with a second transition in every full period.
\end{proof}

\begin{proposition}[Energy of one periodic cell]\label{prop:phase-periodic-cell-energy}
Let $\mathcal E(a)=\int_0^{T(a)}((f_a')^2/2+W(f_a))dt$. Then
\[
 \mathcal E(a)=4\int_0^a\sqrt{2(W(s)-W(a))}\,ds+W(a)T(a),
 \qquad \mathcal E'(a)=W(a)T'(a)>0.
\]
Its endpoint limits are $\pi/2$ at $a=0$ and $2c_W$ at $a=1$.
\end{proposition}
\begin{proof}
The Hamiltonian says that the energy density is $(f_a')^2+W(a)$. Substitution of the profile value on each quarter-period therefore gives the first formula. Write the first term as $J(a)$. The integrand vanishes at its upper endpoint, and its parameter derivative has only an integrable inverse-square-root singularity there. Differentiation, justified either by the substitution $s=a\sin\theta$ or by truncation followed by that bound, yields
\[
 J'(a)=-4W'(a)\int_0^a\frac{ds}{\sqrt{2(W(s)-W(a))}}
       =-W'(a)T(a).
\]
The derivative of $J+WT$ is consequently $WT'$, proving strict increase. At zero $J\to0$, $W(a)\to1/4$, and $T\to2\pi$. At one the integrand of $J$, extended by zero from $[0,a]$ to $[0,1]$, is bounded by $\sqrt{2W(s)}$ and converges to it. Hence $J\to4\int_0^1\sqrt{2W}=2c_W$. Finally $W(a)=(1-a^2)^2/4$ times the logarithmic period bound tends to zero. This proves the second limit without confusing a full period with one transition.
\end{proof}

\begin{example}[All nonconstant phases on a circle]\label{ex:phase-circle-periods}
On a circle of length $\ell$, the equation $\varepsilon^2u''=u^3-u$ has a nonconstant solution with exactly $m$ full oscillations precisely when
\[
 m\in\mathbb N,\qquad 2\pi m\varepsilon<\ell.
\]
Its amplitude is the unique $a\in(0,1)$ solving $T(a)=\ell/(m\varepsilon)$, and it is unique up to translation and reversal. At equality no nonconstant solution of that oscillation number exists.
\end{example}
\begin{proof}
Lift a solution periodically to the line and rescale its independent variable by $\varepsilon$. The bounded classification excludes heteroclinics and leaves exactly the periodic profiles. Their least period in the physical coordinate is $\varepsilon T(a)$, so fitting $m$ complete periods is equivalent to the displayed equation. Proposition \ref{prop:phase-period-function} maps $(0,1)$ strictly increasingly onto $(2\pi,\infty)$, proving existence and uniqueness exactly under the strict inequality. At the endpoint the amplitude would be zero, giving the constant zero rather than a new periodic profile.
\end{proof}

\begin{example}[An explicit negative eigenfunction in every periodic cell]\label{ex:phase-periodic-negative-mode}
For a periodic profile of amplitude $a$, set
\[
 d_a=\sqrt{1+3(1-a^2)^2},\qquad
 b_a=\frac{2+d_a}{3},\qquad v_a(t)=b_a-f_a(t)^2.
\]
Then $v_a>0$ is periodic and
\[
 \left(-\frac{d^2}{dt^2}+3f_a^2-1\right)v_a
   =\lambda_a v_a,\qquad \lambda_a=1-d_a<0.
\]
On a circle comprising one or any integer number of periods, $\lambda_a$ is the least Rayleigh quotient, attained only by multiples of $v_a$. Moreover
$\lambda_a=-\tfrac32(1-a^2)^2+O((1-a^2)^4)$ as $a\uparrow1$.
\end{example}
\begin{proof}
Put $B=a^2(2-a^2)$. The Hamiltonian gives
$(f_a')^2=f_a^4/2-f_a^2+B/2$, whence
$(f_a^2)''=3f_a^4-4f_a^2+B$.
Consequently, for a constant $b$,
\[
 \left(-D^2+3f_a^2-1\right)(f_a^2-b)
       =(3-3b)f_a^2+b-B.
\]
It is an eigenfunction with eigenvalue $3-3b$ precisely when
$3b^2-4b+B=0$. The larger root is $b_a$, giving $\lambda_a=1-d_a$. Since $a<1$, $d_a>1$ and $b_a>1>a^2$, so $v_a$ is strictly positive.

For any periodic $H^1$ test $\phi$, expansion and integration by parts around the circle give
\[
 Q_f(\phi)-\lambda_a\int\phi^2
   =\int v_a^2\left|(\phi/v_a)'\right|^2.
\]
There is no endpoint term because the functions and coefficients are periodic. This proves the least-quotient assertion and its equality characterization without a spectral completeness theorem. Taylor expansion of $\sqrt{1+3(1-a^2)^2}$ gives the final estimate. The negative direction becomes small in a long cell, but it remains negative at every finite amplitude.
\end{proof}

\begin{proposition}[Periodic profiles have infinite index on the line]\label{prop:phase-periodic-infinite-index}
Every nonconstant periodic entire profile has infinite Morse index for compactly supported variations. Its extension $u(y,t)=f_a(t)$ to every $\mathbb R^{d+1}$, $d\geq1$, also has infinite index. Here infinite index means that the quadratic form is negative definite on compactly supported test subspaces of arbitrarily large finite dimension.
\end{proposition}
\begin{proof}
Example \ref{ex:phase-periodic-negative-mode} supplies a positive bounded periodic eigenfunction $v_a$ with negative eigenvalue $\lambda_a$. For a nonzero compact smooth cutoff $\chi$, the same integration by parts on the line gives
\[
 Q_f(v_a\chi(t/R))
 =\lambda_a\int v_a^2\chi(t/R)^2+
   R^{-2}\int v_a^2\chi'(t/R)^2<0
\]
for large $R$: the first term is bounded above by a negative constant times $R$, while the second is $O(R^{-1})$. Thus there is a compactly supported smooth negative test $\phi$. This also gives a concrete alternative to the endpoint-pairing proof of instability in Lemma \ref{lem:phase-stable-ode}.

Translate $\phi$ by sufficiently separated integer multiples of $T(a)$. The potential is periodic, each form value is the same negative number, and disjoint supports eliminate all cross terms. Their span has negative definite form, with arbitrarily large dimension. In higher dimension test with $\phi(t)\eta(y/R)$ for a fixed nonzero compact smooth $\eta$. Dividing by $\int\eta(y/R)^2$ leaves $Q_f(\phi)+R^{-2}C_\eta\int\phi^2$, which is negative for large $R$. Separate copies by integer periods in $t$ to obtain the same arbitrarily large negative subspaces. This strengthens the instability conclusion without needing a full Floquet-spectrum calculation.
\end{proof}
